
\section{Operational Model}


\begin{figure}[H]
\centering
\begin{tikzpicture}[every node/.style={figlabel}, box/.style={figbox, minimum width=3.1cm, minimum height=.9cm, align=center}]
\path[figpanel] (-1.25,-3.15) rectangle (9.25,3.0);
\node[figtitle, anchor=west] at (-0.8,2.55) {SCOPED RESOLUTION};
\node[box, fill=BitBlue!8, draw=BitBlue!70, line width=1.0pt] (root) at (4,1.85) {root\\\statefile{STATE.md}};
\node[box, fill=BitGray, draw=BitLine, text=BitSlate] (pa) at (0.2,0.35) {project-a\\\statefile{STATE.md}};
\node[box, fill=BitTeal!8, draw=BitTeal!80, line width=1.0pt] (sub) at (4,0.35) {subsystem\\\statefile{STATE.md}};
\node[box, fill=BitGray, draw=BitLine, text=BitSlate] (pb) at (7.8,0.35) {project-b\\\statefile{STATE.md}};
\node[box, fill=BitGold!14, draw=BitGold!80!black, line width=1.0pt] (task) at (4,-1.15) {task-42\\\statefile{STATE.md}};
\draw[BitLine,line width=1.6pt] (root) -- (pa);
\draw[BitLine,line width=1.6pt] (root) -- (pb);
\draw[BitBlue,line width=2.0pt] (root) -- (sub) -- (task);
\node[neutralbox, minimum width=8.4cm, minimum height=.92cm, align=center] at (4,-2.35) {Load the smallest sufficient state chain. Ignore unrelated branches unless the task touches them.};
\end{tikzpicture}
\caption{Scoped resolution follows the active branch and leaves unrelated repository state cold.}
\label{fig:scope-resolution}
\end{figure}


A Continuity workspace can be modeled as a tuple
\begin{equation}
\label{eq:workspace}
C = \langle M, S, T, L, H, A, G, V, \Omega \rangle,
\end{equation}
where $M$ is the manifest set, $S$ the current-state set, $T$ the task queue, $L$ learned memory, $H$ semantic history, $A$ generated artifacts, $G$ version history, $V$ validators and enforcement mechanisms, and $\Omega$ the active agent set.

For a path $p$ and context budget $B$, Continuity does not load every available file. It loads the smallest relevant slice:
\begin{equation}
\label{eq:load}
load(p,B)=\operatorname{select}(M\cup S\cup T\cup L\mid scope(p)) \quad \text{s.t.} \quad tokens(load)\leq B.
\end{equation}
The function is intentionally abstract. Its purpose is normative, not implementation-specific: keep the live context small, task-specific, and explicit.

\subsection{Reconciliation invariant}

Every action that changes repository reality risks invalidating recorded state. Continuity therefore requires a reconciliation step after meaningful work:
\begin{equation}
\label{eq:reconcile}
state(C_{after}) \approx reality(R_{after}).
\end{equation}
The approximation symbol is deliberate. Natural-language state files cannot perfectly encode repository truth. They can, however, remain close enough to prevent stale planning and repeated errors.


\begin{figure}[H]
\centering
\begin{tikzpicture}[every node/.style={figlabel}, file/.style={figbox, minimum width=3.55cm, minimum height=1.28cm, align=center}]
\path[figpanel] (-1.05,-2.35) rectangle (10.2,2.55);
\node[figtitle, anchor=west] at (-0.6,2.1) {RECONCILIATION};
\node[file, fill=BitTeal!8, draw=BitTeal!75] (old) at (0.9,0) {\statefile{STATE.md}\\\figsmall API not deployed};
\node[draw, circle, fill=BitGold!18, draw=BitGold!80!black, minimum size=1.28cm, font=\sffamily\small] (change) at (4.55,0) {change};
\node[file, fill=BitTeal!8, draw=BitTeal!75] (new) at (8.3,1.05) {\statefile{STATE.md}\\\figsmall API deployed};
\node[file, fill=BitRose!8, draw=BitRose!72] (hist) at (8.3,-1.15) {\file{snapshots/}\\\figsmall pre-deploy state};
\draw[figarrow] (old) -- (change);
\draw[figflow] (change) -- node[above, figsmall, fill=white, inner sep=1.5pt] {overwrite present} (new);
\draw[figdanger] (change) -- node[below, figsmall, fill=white, inner sep=1.5pt] {append past} (hist);
\end{tikzpicture}
\caption{The live state updates in place, while the old state moves to append-only history.}
\label{fig:reconciliation-loop}
\end{figure}


\subsection{Safety inheritance}

Local context should refine parent guidance, not weaken it. If a repository-level file forbids destructive commands without approval, a subdirectory file must not silently remove that rule. Continuity expresses that constraint as
\begin{equation}
\label{eq:safety}
safety(child) \supseteq safety(parent).
\end{equation}

\subsection{Evidence requirement}

Completion without evidence is only an assertion. Continuity therefore treats a task as complete only when a reviewable artifact exists:
\begin{equation}
\label{eq:evidence}
status(task)=complete \Rightarrow \exists e\in(A\cup G\cup H): supports(e,task).
\end{equation}
A passing test log, an artifact bundle, a review record, or a commit can satisfy this condition.
